\documentclass{beamer}

\usetheme{Madrid} % Sets details for virtually all stylistic choices in presentation

\usepackage{amsmath,multicol,tikz,url,mathrsfs}
\usepackage[T1]{fontenc} % Output accented letters in correct font (i.e. consistently with font setting)
\usepackage[utf8]{inputenc} % Input accented letters directly from keyboard
\usepackage[absolute,overlay]{textpos} % Textblock for position specification


\definecolor{MDB}{RGB}{33, 49, 77} % Mines Dark Blue
\definecolor{MLB}{RGB}{146, 162, 189} % Mines Light Blue
\definecolor{MA}{RGB}{210, 73, 42} % Mines Accent
\definecolor{MB}{RGB}{38, 63, 106} % Mines Blue
\definecolor{MPB}{RGB}{206, 213, 221} % Mines Pale Blue
\definecolor{MCG}{RGB}{178, 180, 179} % Mines Cool Gray
\definecolor{MM}{RGB}{139, 141, 142} % Mines Metallic


\setbeamercolor{palette primary}{bg=MDB,fg=white}
\setbeamercolor{palette secondary}{bg=MLB,fg=white}
\setbeamercolor{palette tertiary}{bg=MB,fg=white}
\setbeamercolor{palette quaternary}{bg=MLG,fg=white}
\setbeamercolor{structure}{fg=MB} % itemize, enumerate, etc
\setbeamercolor{section in toc}{fg=MB} % TOC sections
\setbeamercolor{block title alerted}{fg=white,bg=MA} % alert blocks
\setbeamercolor{block body alerted}{fg=MDB,bg=MA!10!white}
\AtBeginEnvironment{alertblock}{\setbeamercolor{item projected}{bg=MA}}
\setbeamercolor{block title example}{fg=white,bg=MM} % example blocks
\setbeamercolor{block body example}{fg=MDB,bg=MM!10!white,use=MM}
\AtBeginEnvironment{exampleblock}{\setbeamercolor{item projected}{bg=MCG}}

%Information to be included in the title page:
\title[\LaTeX Tutorial]{AMS Bootcamp: LaTeX Tutorial}
\author{Justin Garrish}
\institute[CSM]{Colorado School of Mines}
\date[17 Aug 2020]{17 August 2020}



\begin{document}
\setbeamertemplate{caption}{\raggedright\insertcaption\par}
\AtBeginSubsection[]
{
  \begin{frame}<beamer>{Outline}
    \tableofcontents[currentsubsection]
  \end{frame}
}

\frame{\titlepage}

\begin{frame}
\frametitle{Outline}
\tableofcontents
\end{frame}

\section{Introduction to LaTeX} %------------------------- Introduction -------------------------%
\begin{frame}{What is latex?}
	\begin{itemize}
		\item Latex is a stable dispersion of polymer microparticles in an aqueous medium.
		\item It is found in nature, but synthetic latexes can be made by polymerizing a monomer such as styrene that has been emulsified with surfactants.
		\item Latex as found in nature is a milky fluid found in 10\% of all flowering plants.
		\item It is a complex emulsion consisting of 
			\begin{itemize}
				\item proteins
				\item alkaloids
				\item starches
				\item sugars
				\item oils\dots
			\end{itemize}
	\end{itemize}
	\only<2>{
	\begin{textblock*}{8cm}(2.5cm,2.5cm) % {block width} (coords)
		\includegraphics[width=8cm]{facepalm.png}
	\end{textblock*}
	}
\end{frame}

\begin{frame}{\LaTeX and \TeX}
	\begin{multicols}{2}
		\begin{itemize}
			\item<1-> TeX is a typesetting system for  writing developed in 1978 by Donald Knuth.
			\item<2-> The letters T, E, X in TeX come from the Greek letters tau, epsilon, and chi.
			\item<3-> LaTeX, a document preparation system containing TeX macros for ease of access, was developed by Leslie Lamport in the 1980s (initial release 1984).
		\end{itemize}
	\columnbreak
		\begin{figure}
			\includegraphics[width = 2.75cm]{don_knuth.png}\ \includegraphics[width = 2.75cm]{les_lamport.png}
			\caption{Donald Knuth (left) and Leslie Lamport (right)}
		\end{figure}
	\end{multicols}
\end{frame}

\begin{frame}{Learning Outcomes}
	By the end of this tutorial, participants will be able to\dots
	\begin{itemize}
		\item Write a simple LaTeX article with conscious awareness of preamble and document content.
		\item Typeset documents containing thoughtfully formatted mathematical expressions/statements.
		\item Utilize common commands and environments relevant to scientific and mathematical communication.
	\end{itemize}
\end{frame}

\begin{frame}{Uses of LaTeX}
	\begin{itemize}
		\item<1-> LaTeX allows users to typeset scientific and technical documents, including
			\begin{itemize}
				\item Journal articles
				\item Presentations
				\item Textbooks
				\item Posters
				\item Course notes (Student\footnote{This typically takes some familiarity with LaTeX and planning in advance to be effective.  Can discuss this afternoon!} or Instructor)
			\end{itemize}
		\item<2-> Today, LaTeX-based equation editors are integrated into several common tools utilized in academia:
			\begin{itemize}
				\item Word processors (e.g. MS Word, Apple Pages)
				\item Presentation tools (e.g. MS PowerPoint, Apple Keynote)
				\item Learning management and testing tools (e.g. Canvas, GradeScope)
			\end{itemize}
	\end{itemize}
\end{frame}

\begin{frame}{Features}
	\begin{itemize}
		\item High-quality typesetting of complicated mathematical formulae
		\item Clean, sensible default layouts to begin writing with minimal effort
		\item Fine-tuned control over stylistic details available as needed or desired
		\item Characters from several alphabets and languages
		\item Dedicated functionality for organizing documents and presenting information
	\end{itemize}
	\vspace{0.5cm}

	\begin{multicols}{3}
		\begin{figure}
			\resizebox{0.275\textwidth}{!}{
			\begin{tikzpicture}
				\draw[thick] (0,0) -- (5,0);
				\draw[thick] (5,0) -- (2.5,4.330127018922193);
				\draw[thick] (2.5,4.330127018922193) -- (0,0);
				\fill[black] (0,0) -- (5,0) -- (2.5,4.330127018922193); \pause
				\draw[thick] (2.5,0) -- (1.25,2.16506);
				\draw[thick] (1.25,2.16505) -- (3.75,2.16505);
				\draw[thick] (3.75,2.16505) -- (2.5,0);
				\fill[white] (2.5,0) -- (1.25,2.16506) -- (3.75,2.16505); \pause
				\draw[thick] (1.25,0) -- (1.875,1.082532);
				\draw[thick] (1.875,1.082532) -- (0.625,1.082532);
				\draw[thick] (0.625,1.082532) -- (1.25,0);
				\fill[white] (1.25,0) -- (1.875,1.082532) -- (0.625,1.082532);
				\draw[thick] (3.75,0) -- (4.375,1.082532);
				\draw[thick] (4.375,1.082532) -- (3.125,1.082532);
				\draw[thick] (3.125,1.082532) -- (3.75,0);
				\fill[white] (3.75,0) -- (4.375,1.082532) -- (3.125,1.082532);
				\draw[thick] (2.5,2.16506) -- (3.125,3.247595);
				\draw[thick] (3.125,3.247595) -- (1.875,3.247595);
				\draw[thick] (1.875,3.247595) -- (2.5,2.16506);
				\fill[white] (2.5,2.16506) -- (3.125,3.247595) -- (1.875,3.247595); \pause
				\fill[white] (0.3125,0.541266) -- (0.625,0) -- (0.9375,0.541266);
				\fill[white] (0.9375,1.623798) -- (1.25,1.082532) -- (1.5625,1.623798);
				\fill[white] (1.5625,0.541266) -- (1.875,0) -- (2.1875,0.541266);
				\fill[white] (1.5625,2.706329) -- (1.875,2.165064) -- (2.1875,2.706329);
				\fill[white] (2.1875,3.788861) -- (2.5,3.247595) -- (2.8125,3.788861);
				\fill[white] (2.8125,2.706329) -- (3.125,2.165064) -- (3.4375,2.706329);
				\fill[white] (2.8125,0.541266) -- (3.125,0) -- (3.4375,0.541266);
				\fill[white] (3.4375,1.623798) -- (3.75,1.082532) -- (4.0625,1.623798);
				\fill[white] (4.0625,0.541266) -- (4.375,0) -- (4.6875,0.541266);
			\end{tikzpicture}
			}
		\caption{Drawn in TikZ}
		\end{figure}
	\columnbreak
	Aligned equations\vspace{-1cm}
		\begin{align*}
			\dot x\ &=\ \mu\left(x-\frac13x^3-y\right)\\
			\dot y\ &=\ \dot y\frac1{\mu}x
		\end{align*}
	\columnbreak
		\begin{figure}
			\centering
			\includegraphics[width = 0.2\textwidth]{turbulent.png}
			\caption{Local .png image} %%%%% Include citation
		\end{figure}
	\end{multicols}
\end{frame}

\section{LaTeX Basics} %------------------------- Basics -------------------------%

\subsection{Hello, World!}

\begin{frame}{LaTeX Commands}
	\begin{itemize}
		\item Indicated by a backslash ``\textbackslash'' preceding the command name
		\item Processed by LaTeX as a set of instructions
		\item General structure:
			\begin{semiverbatim}
				\\commandname[<optional arguments>]\{<mandatory arguments>\}
			\end{semiverbatim}
	\end{itemize}
	\onslide<4->{
	\begin{block}{Note}
		Not all commands have optional arguments.
	\end{block}
	}
\end{frame}

\begin{frame}{Document Class}
	\begin{semiverbatim}
		\\documentclass[<options>]\{<class>\}
	\end{semiverbatim}
	\begin{itemize}
		\item<1-> The document class sets the default layout and format options 
		\item<2-> Some common classes are
			\begin{itemize}
				\item article
				\item report
				\item beamer
				\item book
				\item minimal
			\end{itemize}
		\item<3-> The available options depend on the choice of class
	\end{itemize}
	\only<4>{
	\begin{exampleblock}{Example}
	Journal article with A4 paper size and 12-point font
		\begin{semiverbatim}
			\\documentclass[a4paper,12pt]\{article\}
		\end{semiverbatim}
	\end{exampleblock}
	}
\end{frame}

\begin{frame}[fragile]{Hello, World!}
	\begin{semiverbatim}
		\\documentclass\{article\}
		
		\\begin\{document\}
		
 			Hello, World!
		
		\\end\{document\}
	\end{semiverbatim}
\end{frame}

\begin{frame}{General Layout}
	The \texttt{Hello, World!} program illustrates the two primary components of a LaTeX document:
		\begin{itemize}
			\item<1-> Preamble
				\begin{itemize}
					\item Document properties (including required \texttt{\textbackslash documentclass} \textit{command})
					\item Packages and libraries
					\item Style and format settings (whole document)
					\item Title-page information
					\item Custom commands
				\end{itemize}
			\item<2-> Document
				\begin{itemize}
					\item All writing and displayed content contained within \texttt{document} \textit{environment}
					\item Many settings may be changed within the document part (e.g. for local formatting)
					\item Can be organized into sections, subsections, paragraphs, etc.
				\end{itemize}
		\end{itemize}
\end{frame}

\subsection{Commands and Environments}

\begin{frame}{Common Commands}
\framesubtitle{Load Packages}
	\textbf{Preamble}
	\begin{semiverbatim}
		\\usepackage[<options>]\{<packages>\}
	\end{semiverbatim}
	\begin{itemize}
		\item Additional commands, options, and environments available
		\item Multiple packages can be listed within the same pair of braces (with same/no options declared)
		\item Packages can inherit options from the \texttt{\textbackslash documentclass} command
	\end{itemize}
	\onslide<2>{
	\begin{exampleblock}{Example}
		Loads packages for mathematics typesetting, enumeration and multiple-column formatting, and graphical tools:
			\begin{semiverbatim}
				\\usepackage\{amsmath, amsthm, amssymb, mathtools, enumitem, multicols, tikz, pgfplots\}
			\end{semiverbatim}
	\end{exampleblock}
	}
\end{frame}

\begin{frame}{Common Commands}
\framesubtitle{Title Page}
	\textbf{Preamble}
	\begin{semiverbatim}
		\\title\{<Document Title>\}
		
		\\author\{<Name>\}
		
		\\date\{<Date>\}
	\end{semiverbatim}
	\textbf{Document}
	\begin{semiverbatim}
		\\maketitle
	\end{semiverbatim}
\end{frame}

\begin{frame}{Common Commands}
\framesubtitle{Sectioning}
	\textbf{Document}
	\begin{semiverbatim}
		\\section\{<Section Title>\}
		
		\\subsection\{<Subsection Title>\}
		
		\\subsubsection\{<Subsubsection Title>\}
	\end{semiverbatim}
	\begin{itemize}
		\item Sections are numbered in the order that they appear within the \texttt{document} environment by default.
		\item An asterisk ``*'' between the command name and braces suppresses the numbering for that section.
	\end{itemize}
	\onslide<2>{
	\begin{exampleblock}{Example}
			\begin{semiverbatim}
				\\section*\{Cheesy Maths Jokes\}
				
				\\subsection*\{Infinitely Many Mathematicians\\dots\}
			\end{semiverbatim}
	\end{exampleblock}
	}
\end{frame}

\begin{frame}{Common Commands}
\framesubtitle{Environments}
	\textbf{Document}
	\begin{semiverbatim}
		\\begin\{<environment>\}
		
		 <Content>
			
		\\end\{<environment>\}
	\end{semiverbatim}
	\onslide<2->{
	\begin{itemize}
		\item Environments alter the typesetting rules
		\item Can be nested when multiple environments are in play
		\item The \texttt{document} environment is necessary at the outermost level of each LaTeX file
	\end{itemize}
	More to come\dots
	}
\end{frame}

\begin{frame}{Common Commands}
\framesubtitle{Miscellany}
	\begin{itemize}
		\item<1-> \texttt{\textbackslash textbf\{\}}, \texttt{\textbackslash textit\{\}}, \texttt{\textbackslash underline\{\}}
		\item<2-> \texttt{\textbackslash textrm\{\}}, \texttt{\textbackslash textsf\{\}}, \texttt{\textbackslash texttt\{\}}, \texttt{\textbackslash textsc\{\}}, \texttt{\textbackslash emph\{\}}
		\item<3-> \texttt{\textbackslash bibliography\{\}}, \texttt{\textbackslash cite\{\}}
		\item<4-> \texttt{\textbackslash label\{\}}, \texttt{\textbackslash ref\{\}}, \texttt{\textbackslash pageref\{\}}, \texttt{\textbackslash footnote\{\}}
		\item<5-> \texttt{\textbackslash tiny}, \texttt{\textbackslash scriptsize}, \texttt{\textbackslash footnotesize}, \texttt{\textbackslash small}, \texttt{\textbackslash large}, \texttt{\textbackslash huge}
		\item<6-> \texttt{\textbackslash centering}, \texttt{\textbackslash raggedright}, \texttt{\textbackslash raggedleft}
		\item<7-> \texttt{\textbackslash dots}, \texttt{\textbackslash cdots}, \texttt{\textbackslash vdots}, \texttt{\textbackslash hdots}, \texttt{\textbackslash ddots}
		\item<8-> \texttt{\textbackslash newline}, \texttt{\textbackslash newpage}, \texttt{\textbackslash vspace\{\}}, \texttt{\textbackslash hspace\{\}}, \texttt{\textbackslash rule\{\}\{\}}
	\end{itemize}
\end{frame}

\begin{frame}{Common Environments}
\framesubtitle{Listing}
	\begin{itemize}
		\item<1-> ``Listing environments'' provide various formats for listed content
		\item<2-> Items in the list are distinguished by the \texttt{\textbackslash item} command.
			\begin{itemize}
				\item \texttt{itemize} -- bulleted list
				\item \texttt{enumerate} -- ordered list\footnote{With package \texttt{enumitem}, the list appearance may be easily customized, e.g. (a), (b), (c) instead of 1., 2., 3.}
				\item \texttt{description} -- Description list; requires command \texttt{\textbackslash item[<description>]}, with \texttt{<description>} used instead of a bullet or number.
			\end{itemize}
	\end{itemize}
	\onslide<3->{
	\begin{exampleblock}{Example}
		\begin{semiverbatim}
			\\begin\{description\}
			
			\\item[First] The first thing
				
			\\item[Second] A second thing
				
			\\end\{description\}
		\end{semiverbatim}
	\end{exampleblock}
	}
\end{frame}

\begin{frame}{Common Environments}
\framesubtitle{Floating Bodies}
	A \textbf{float} is a container for content in a document that cannot be broken over a page (as opposed to, say, freely typed text).
	\begin{semiverbatim}
		\\begin\{<float>\}[placement]
		
		 <Content>
			
		\\end\{<float>\}
	\end{semiverbatim}
	\begin{itemize}
		\item<2-> Common floats:
			\begin{itemize}
				\item \texttt{table}
				\item \texttt{figure}
			\end{itemize}
		\item<3-> Placement options
			\begin{itemize}
				\item t -- ``top''
				\item h -- ``here''
				\item b -- ``bottom''
				\item p -- ``separate page''
				\item ! -- ``override \LaTeX parameters''
			\end{itemize}
	\end{itemize}
\end{frame}

\begin{frame}{Common Environments}
\framesubtitle{Float Example}
	\setbeamertemplate{caption}[default]
	\begin{figure}[h]
		\includegraphics[width=0.3\textwidth]{math_universal.png}
		\label{fig:math_universal}
		\caption{Re-creation of South Park's iconic ``Math is universal'' poster. https://andyborne.com/math/freestuff.html}
	\end{figure}
\end{frame}

\begin{frame}{Common Environments}
\framesubtitle{Float Example}
	\begin{exampleblock}{Math is universal}
		\begin{semiverbatim}
			\\begin\{figure\}[h]
			
			\\includegraphics[width=0.3\\textwidth]\{math\_universal.png\}
			
			\\label\{fig:math\_universal\}
			
			\\caption\{Re-creation of South Park's iconic ``Math is universal'' poster. https://andyborne.com/math/freestuff.html\}

			\\end\{figure\}
		\end{semiverbatim}
	\end{exampleblock}
	\begin{alertblock}{Notes}
		\begin{itemize}
			\item The commands \texttt{\textbackslash includegraphics}, \texttt{\textbackslash caption}, and \texttt{\textbackslash label} belong \emph{inside} the \texttt{figure} (or other \texttt{<float>}) environment
			\item The marker used for \texttt{\textbackslash label\{<marker>\}} is the same one used to reference this figure elsewhere in the document with \texttt{\textbackslash ref\{<marker>\}}
		\end{itemize}
	\end{alertblock}
\end{frame}

\begin{frame}{Common Environments}
\framesubtitle{Tables}
	\begin{semiverbatim}
		\\begin\{tabular\}[<position>]\{<columns>\}
		
		 <Content>
			
		\\end\{tabular\}
	\end{semiverbatim}
	\begin{itemize}
		\item<2-> The \texttt{tabular} environment produces a sequence of rows of items, aligned vertically (in columns).
		\item<3-> Position does not need to be specified by default, and can also be inherited if the \texttt{tabular} environment is nested within a \texttt{table} (float) environment.
		\item<4-> The \texttt{<columns>} option conveys how many columns the table contains and the justification (alignment) of each:
			\begin{itemize}
				\item l -- left-justified
				\item c -- center-justified
				\item r -- right-justified
			\end{itemize}
	\end{itemize}
\end{frame}

\begin{frame}{Common Environments}
	\framesubtitle{Table Example}
	\begin{exampleblock}{Table of Greek letters}
		\begin{semiverbatim}
			\\begin\{table\}[t]
			
			\\caption\{Some Greek Letters\}
			
				\\begin\{tabular\}\{| c | c | c |\}\\hline
				
				Letter \& Lowercase \& Uppercase \\\\\ \\hline
				
				alpha \& \$\\alpha\$ \& \$\\mathrm A\$ \\\\\ \\hline
				
				beta \& \$\\beta\$ \& \$\\mathrm B\$ \\\\\ \\hline
				
				gamma \& \$\\gamma\$ \& \$\\Gamma\$ \\\\\ \\hline
				
				\\end\{tabular\}
				
			\\label\{tbl:greek\}

			\\end\{table\}
		\end{semiverbatim}
	\end{exampleblock}
\end{frame}

\begin{frame}{Common Environments}{Table Example}
	\begin{table}[t]
		\caption{Some Greek Letters}
		\begin{tabular}{| c | c | c |}\hline
			Letter & Lowercase & Uppercase\\ \hline
			alpha & $\alpha$ & $\mathrm A$ \\ \hline
			beta & $\beta$ & $\mathrm B$ \\ \hline
			gamma & $\gamma$ & $\Gamma$ \\ \hline
		\end{tabular}
		\label{tbl:greek}
	\end{table}
	\onslide<2->{
	\begin{alertblock}{Notes}
		\begin{itemize}
			\item Vertical bars within the \texttt{<columns>} option add vertical lines at corresponding boarders to columns
			\item The command \texttt{\textbackslash\textbackslash} moves to the next row of the table
			\item The command \texttt{\textbackslash hline} adds a horizontal line at the boarder above the following row
		\end{itemize}
	\end{alertblock}
	}
\end{frame}

\begin{frame}{Common Environments}
\framesubtitle{Text Environments}
Text environments create a block of text with modified typesetting rules for various circumstances:
	\begin{itemize}
		\item \texttt{quote} -- Indented quote block
		\item \texttt{quotation} -- Like quote, but with indented paragraphs
		\item \texttt{verse} -- Like quotation, but indentation inverted
		\item \texttt{verbatim} -- Monospaced text with literal characters (good for replicating computer code in LaTeX)
		\item \texttt{verbatim*} -- Like verbatim, but with spaces as \texttt{\char32}
	\end{itemize}
\end{frame}

\subsection{Bibliographies}

\begin{frame}{BibTeX}
\framesubtitle{Background}
	\begin{itemize}
		\item BibTeX is a reference management and typesetting tool designed for use with LaTeX.
		\item<2-> Bibliographic information for each source is written in plain text and saved with the \texttt{.bib} extension.
		\item<3-> Editors (e.g. BibDesk, EndNote) exist to provide an interface beyond plain text
		\item<4-> References are formatted automatically based on the \texttt{\textbackslash bibliographystyle\{<style>\}} (preamble) command.
		\item<5-> The command \texttt{\textbackslash bibliography\{<filename>\}} places the bibliography in the document.
		\item<6-> Each source in BibTeX has an associated reference key (user-determined) for use in citations.
			\begin{itemize}
				\item The command \texttt{\textbackslash cite\{<key>\}} will place an in-line citation to the corresponding reference.
				\item Only sources that are cited in line will be added to the displayed reference list.
			\end{itemize}
	\end{itemize}
\end{frame}

\begin{frame}[fragile]{BibTeX}
\framesubtitle{Example BibTeX Code}
The following would be placed in the \texttt{.bib} file corresponding to the LaTeX document:
	\begin{exampleblock}{}
		\begin{semiverbatim}
			@article\{Man2002, \% cite key
			title=\{The oral glucose minimal model: Estimation of insulin sensitivity from a meal test\},
			author=\{C. D. Man and A. Caumo and C. Cobelli\},
			journal=\{IEEE Transactions on Biomedical Engineering\},
			year=\{2002\},
			volume=\{49\},
			pages=\{419-429\}
			\}
		\end{semiverbatim}
	\end{exampleblock}
\end{frame}

\section{Typesetting Mathematics} %------------------------- Typesetting Mathematics -------------------------%

\subsection{Inline Mathematics}

\begin{frame}{Math Mode}
In order for text within a paragraph to be typeset as mathematics, a special environment is used: \texttt{\$\dots\$} or \texttt{\textbackslash(\dots\textbackslash)}.

	\begin{exampleblock}{Example}
		Euler's Identity: \texttt{\$e\^{}\{i\textbackslash pi\}+1=0\$}
	\end{exampleblock}	\begin{itemize}
		\item<2-> \textbf{Math mode} refers to the designation of special typesetting rules in effect within certain environments.
		\item<3-> Normal text is not processed as such within math mode (words are processed like strings of successive variables, and spaces are skipped, for instance).
		\item<4-> Math mode offers access to a wide assortment of commands relevant to mathematics, and to scientific and technical writing involving specialized notation
	\end{itemize}
\end{frame}

\begin{frame}{Mathematical Commands and Operators}
\framesubtitle{A very small sample}
	\begin{center}
		\begin{tabular}{| l | c || l | c |}\hline
			Command & Example & Command & Example\\ \hline
			\texttt{\textbackslash rightarrow} & $\rightarrow$ & \texttt{\textbackslash Rightarrow} & $\Rightarrow$ \\ \hline
			\texttt{\textbackslash rightharpoonup} & $\rightharpoonup$ & \texttt{\textbackslash frac\{a\}\{b\}} & $\frac ab$  \\ \hline
			\texttt{\textbackslash in} & $\in$ & \texttt{\textbackslash ni} & $\ni$ \\ \hline
			\texttt{\textbackslash sum\_\{j=1\}\^{}n} & $\sum_{j=1}^n$ & \texttt{\textbackslash prod\_\{j=1\}\^{}n} & $\prod_{j=1}^n$ \\ \hline
			\texttt{\textbackslash partial} & $\partial$ & \texttt{\textbackslash int\_a\^{}b} & $\int_a^b$ \\ \hline
			\texttt{\textbackslash iint\_a\^{}b} & $\iint_a^b$ & \texttt{\textbackslash oint\_a\^{}b} & $\oint_a^b$ \\ \hline
		\end{tabular}
		
		\vspace{0.5cm}
		
		\onslide<2->{
		\begin{tabular}{| l | c || l | c |}\hline
			Operator & Example & Operator & Example\\ \hline
			\texttt{\textbackslash lim\_\{x\textbackslash to\textbackslash infty\}} & $\lim_{x\to\infty}$ & \texttt{\textbackslash max\_\{n\textbackslash le c\}} & $\max_{n\le c}$ \\ \hline
			\texttt{\textbackslash arcsin} & $\arcsin$ & \texttt{\textbackslash cos\^{}2} & $\cos^2$ \\ \hline
			\texttt{\textbackslash exp} & $\exp$ & \texttt{\textbackslash ln} & $\ln$ \\ \hline
		\end{tabular}
		}
	\end{center}
\end{frame}

\subsection{Display Mathematics}

\begin{frame}{Display Mode}
\framesubtitle{Math Environments}
	\begin{itemize}
		\item Several environments exist that typeset mathematical symbols in a larger format than their inline counterparts.
			\begin{itemize}
				\item \texttt{\textbackslash begin\{equation\}\dots\textbackslash end\{equation\}} -- numbered equation (single label for environment, regardless of number of equations)
				\item \texttt{\textbackslash[\dots\textbackslash]} -- equation with numbering suppressed (can use \texttt{equation*} environment, too)
				\item \texttt{\textbackslash begin\{align\}\dots\textbackslash end\{align\}} -- multiple equations (or columns of equations) with vertical alignment specified by ``\&'' (one label per equation)
			\end{itemize}
		\item<2-> Within the environments above, math is set in ``display style'' (as opposed to the ``text style'' of inline math)
		\item<3-> Typically used for the mathematical statements within important results, esp. to be referenced elsewhere
	\end{itemize}
\end{frame}

\begin{frame}{Display Mode}
\framesubtitle{Single-Equation Example}
	\begin{exampleblock}{Code}
		\begin{semiverbatim}
		\textbackslash begin\{equation\}
		
		\textbackslash forall\textbackslash varepsilon>0,\textbackslash,\textbackslash exists\textbackslash delta>0\mbox{ such that }\textbackslash left\textbackslash lvert x-c\textbackslash right\textbackslash rvert>\textbackslash delta\textbackslash Rightarrow\textbackslash left\textbackslash lvert f(x)-L\textbackslash right\textbackslash rvert>\textbackslash varepsilon.
		
		\textbackslash end\{equation\}
		\end{semiverbatim}
	\end{exampleblock}
	
	\begin{equation}
		\forall\varepsilon>0,\,\exists\delta>0\mbox{ such that }\left\lvert x-c\right\rvert<\delta\Rightarrow\left\lvert f(x)-L\right\rvert<\varepsilon.
	\end{equation}
\end{frame}

\begin{frame}[fragile]{Display Mode}
\framesubtitle{Multiple-Equation Example}
	\begin{exampleblock}{Code}
		\begin{semiverbatim}
			\textbackslash begin\{align\}
				\textbackslash{}dot G(t)\&=-\textbackslash{}left[p\_1+X(t)\textbackslash{}right]G(t)+p\_1\textbackslash{}cdot G\_b\&
				    G(0)\&=G\_b,\textbackslash\textbackslash
				\textbackslash{}dot X(t)\&=-p\_2X+p\_3\textbackslash{}left[I(t)-I\_b\textbackslash{}right],\&
				    X(0)\&=0.
			\textbackslash end\{align\}
		\end{semiverbatim}
	\end{exampleblock}
	
	\begin{align}
		\dot G(t)&=-\left[p_1+X(t)\right]G(t)+p_1\cdot G_b,&
			G(0)&=G_b,\\
		\dot X(t)&=-p_2X+p_3\left[I(t)-I_b\right],&
			X(0)&=0.
	\end{align}
	Numbering may be suppressed by using environment \texttt{align*}.
\end{frame}

\begin{frame}{Delimiters}
	\begin{itemize}
		\item Delimiters act logically like parentheses
		\item Can only be used in math mode (though the related symbols may be used in text mode)
		\item Delimiters should ``fit'' around the expressions they encompass: precede with \texttt{\textbackslash left} and \texttt{\textbackslash right} for auto-sizing.
	\end{itemize}\vspace{-0.5cm}
	\onslide<2->{
	\begin{table}[t]
		\caption{\LaTeX{} Delimiters}
		\begin{tabular}{| l | c |}\hline
			Delimiter & Text\\ \hline
			Parentheses & \texttt{( )} \\ \hline
			Brackets & \texttt{[ ]} \\ \hline
			Braces & \texttt{\textbackslash \{ \textbackslash \}} \\ \hline
			Angle Brackets & \texttt{\textbackslash langle \textbackslash rangle}  \\ \hline
			Vertical Lines & \texttt{\textbackslash lvert \textbackslash rvert} \\ \hline
			Double Vertical Lines & \texttt{\textbackslash lVert \textbackslash rVert} \\ \hline
		\end{tabular}
	\end{table}
	}
\end{frame}

\begin{frame}{Math Fonts}
	\begin{table}
		\centering
		\begin{tabular}{| l | c | c | c |}\hline
			Font & Command & Examples & Required Packages \\ \hline\hline
			Blackboard & \texttt{\textbackslash mathbb\{\}} & $\mathbb{N},\mathbb{Z},\mathbb{R},\mathbb{C}$ & \texttt{amssymb}\\ \hline
			Calligraphic & \texttt{\textbackslash mathcal\{\}} & $\mathcal{R},\mathcal{B},\mathcal{P}$ & None\\ \hline
			Roman & \texttt{\textbackslash mathrm\{\}} & $\mathrm{A},\mathrm{B},\mathrm{K}$& None\\ \hline
			Italic & \texttt{\textbackslash mathit\{\}} & $\mathit{123},\mathit{abc},\mathit{ABC}$& None\\ \hline
			Script & \texttt{\textbackslash mathscr\{\}} & $\mathscr{C},\mathscr{L},\mathscr{B}$ & \texttt{mathrsfs}\\ \hline
		\end{tabular}
	\end{table}
\end{frame}

\begin{frame}{Additional Resources}
	\begin{itemize}
		\item \url{https://www.overleaf.com/learn}
		\item \url{https://www.latex-tutorial.com}
		\item \url{https://tex.stackexchange.com}
		\item \url{https://detexify.kirelabs.org/classify.html}
		\item \url{http://ctan.math.washington.edu/tex-archive/macros/latex/required/amsmath/amsldoc.pdf}
		\item Workshop today, 1-3 pm (yay!)
	\end{itemize}
\end{frame}






\end{document}

